\chapter{Kết luận và Hướng Phát triển}

% ============================================================
\section{Kết luận}
\label{sec:conclusion}
% ============================================================

Khóa luận này xuất phát từ một quan sát thực tế: mặc dù greybox fuzzing
là hướng tiếp cận mạnh mẽ nhất hiện nay trong kiểm thử bảo mật tự động,
toàn bộ các công cụ thuộc thế hệ này --- bao gồm Phuzz, webFuzz, và
Witcher --- đều bỏ sót hoàn toàn lớp lỗ hổng Blind SSRF. Nguyên nhân
cốt lõi là vòng lặp fuzzing truyền thống chỉ quan sát hai tín hiệu:
độ phủ mã và ngoại lệ. Blind SSRF không để lại dấu vết trong cả hai
kênh này --- ứng dụng vẫn chạy bình thường, không có exception, không
có dòng code mới, chỉ có một kết nối mạng âm thầm được thực hiện ở
tầng hệ điều hành.

Hướng giải quyết được đề xuất là \textbf{bổ sung kênh tín hiệu
Out-of-Band} vào vòng lặp fuzzing, thay vì cố gắng suy luận từ các
tín hiệu hiện có. Nếu dấu vết duy nhất của Blind SSRF là kết nối mạng
tới địa chỉ ngoài mong đợi, ta đặt một máy chủ lắng nghe tại địa chỉ
đó và đưa tín hiệu callback trở lại vòng lặp như một kênh phản hồi độc
lập.

\subsection{Đóng góp của khóa luận}

Khóa luận đã thiết kế và cài đặt đầy đủ bốn thành phần mới tích hợp
vào nền tảng Phuzz:

\textbf{1. Bộ đột biến SSRF chuyên biệt.} \texttt{SSRFMutator} sinh
payload OOB theo round-robin qua 21 biến thể, tổ chức thành ba nhóm
(IP-based, DNS-based, hostname-based) để bao phủ các cơ chế lọc đầu
vào phổ biến. \texttt{CmdInjectionSSRFMutator} bổ sung 10 biến thể
command injection cho các endpoint có nguy cơ lệnh shell. Thiết kế
nhiều biến thể được chứng minh là cần thiết trong thực tế: hai trong
sáu CVE được kiểm thử --- Canto (CVE-2020-28975) và elFinder
(CVE-2021-32682) --- chỉ phát hiện được qua biến thể đặc thù
(scheme-less bypass và hostname indirection), không phải biến thể
plain IP thông thường.

\textbf{2. Máy chủ OOB lắng nghe.} Container \texttt{phuzz-oob} chạy
Flask HTTP/HTTPS listener và dnsmasq phân giải wildcard DNS
\texttt{*.oob}. Cơ chế ghi log nguyên tử (\texttt{os.rename}) trên
shared-tmpfs đảm bảo fuzzer không bao giờ đọc file chưa hoàn chỉnh.
Thiết kế container độc lập cho phép tái sử dụng OOB server với bất kỳ
ứng dụng mục tiêu nào mà không cần sửa đổi.

\textbf{3. Lớp hook PHP hai tầng.} \texttt{SSRFVulnCheck} kết hợp
tín hiệu in-band (PHP hook ghi nhận lời gọi hàm mạng) và tín hiệu OOB
(callback từ ứng dụng mục tiêu) theo logic AND: chỉ khi cả hai file
bằng chứng cùng tồn tại mới xác nhận lỗ hổng. Module hook
\texttt{08\_ssrf.php} triển khai 23 hook trên 5 nhóm hàm qua UOPZ,
bao phủ từ \texttt{curl\_exec}, \texttt{file\_get\_contents}, các hàm
socket, shell execution, đến WordPress HTTP API wrappers.

\textbf{4. PHP stream wrapper.} \texttt{09\_ssrf\_stream.php} đăng ký
lại built-in wrapper cho scheme \texttt{http} và \texttt{https}, cho
phép interceptsâu language construct \texttt{include} và
\texttt{require\_once} mà UOPZ không thể hook. Lớp này lấp khoảng
trống phát hiện cho lớp lỗ hổng Remote File Inclusion (RFI) --- được
xác nhận qua thực nghiệm với bWAPP và Mutillidae.

\subsection{Kết quả thực nghiệm}

Thực nghiệm trên 14 ca kiểm thử (8 ứng dụng lab với tổng cộng 11
endpoint, 6 CVE thực tế) cho thấy:

\begin{itemize}
    \item \textbf{Tỷ lệ phát hiện: 14/14 (100\%)} --- toàn bộ lỗ hổng
    SSRF trên các mục tiêu đã kiểm thử đều được phát hiện.
    \item \textbf{False positive: 0} --- cơ chế AND logic hai tầng
    đảm bảo mỗi báo cáo đều có bằng chứng mạng thực tế.
    \item \textbf{Phuzz gốc: 0/14} --- xác nhận thực nghiệm rằng tín
    hiệu truyền thống hoàn toàn bất lực trước Blind SSRF.
    \item Năm loại sink function khác nhau được bao phủ:
    \texttt{file\_get\_contents}, \texttt{curl\_exec},
    \texttt{include}/\texttt{require\_once} (qua stream wrapper),
    \texttt{shell\_exec}, và \texttt{getimagesize} --- hai loại cuối là
    sink không trực quan mà hầu hết công cụ khác bỏ sót.
\end{itemize}

Kết quả này khẳng định rằng \textbf{mở rộng kênh tín hiệu phản hồi
là hướng tiếp cận khả thi và hiệu quả} để lấp khoảng trống mà các
greybox fuzzer truyền thống để lại.

% ============================================================
\section{Hạn chế}
\label{sec:limitations}
% ============================================================

Dù đạt kết quả khả quan, hệ thống hiện tại có một số hạn chế cần
nhìn nhận thẳng thắn:

\textbf{Phạm vi ngôn ngữ.} Hệ thống chỉ hỗ trợ PHP. Các ứng dụng
viết bằng Python (Django, Flask), Node.js, hay Java/Spring đều nằm
ngoài tầm bao phủ hiện tại, dù kiến trúc OOB listener và cơ chế
phát hiện hai tầng về nguyên lý có thể áp dụng cho bất kỳ ngôn ngữ
nào.

\textbf{Stored SSRF.} Hệ thống không xử lý được Stored SSRF --- dạng
tấn công trong đó payload được lưu vào database và kích hoạt SSRF ở
một request hoàn toàn khác sau đó. Trong trường hợp này, token từ
request gốc và OOB callback từ request thứ cấp không có cơ chế tự
động liên kết.

\textbf{Thời gian phát hiện trên WordPress.} CVE-2020-24148 mất
khoảng 506 giây để phát hiện, chủ yếu do overhead từ WordPress
bootstrap và quá trình đăng nhập. Với giới hạn thời gian 600 giây,
biên độ an toàn không lớn; các ứng dụng WordPress phức tạp hơn có
thể vượt giới hạn này.

\textbf{Quy mô thực nghiệm.} 14 ca kiểm thử là đủ để kiểm chứng
nguyên lý, nhưng chưa đủ để kết luận thống kê về tỷ lệ phát hiện
tổng quát. Thực nghiệm trên tập lỗ hổng lớn hơn (hàng trăm CVE SSRF)
sẽ cung cấp đánh giá tin cậy hơn.

\textbf{Phụ thuộc vào kết nối mạng từ container.} OOB server phải
tiếp cận được từ container web. Trong môi trường production với tường
lửa egress nghiêm ngặt, container web không thể kết nối ra ngoài và
cơ chế OOB sẽ không hoạt động. Tuy nhiên đây là đặc điểm của môi
trường kiểm thử cục bộ (fuzzing lab), không phải môi trường production.

% ============================================================
\section{Hướng Phát triển Tương lai}
\label{sec:future}
% ============================================================

\textbf{Mở rộng sang ngôn ngữ khác.} Hướng tiếp cận tự nhiên nhất
là xây dựng module hook tương đương cho Python (\texttt{urllib},
\texttt{requests}, \texttt{httpx}), Node.js (\texttt{http},
\texttt{https}, \texttt{fetch}), và Java (\texttt{HttpURLConnection},
\texttt{OkHttpClient}). OOB server và logic phát hiện hai tầng có thể
giữ nguyên; chỉ cần thay module hook theo ngôn ngữ mục tiêu.

\textbf{Xử lý Stored SSRF.} Một hướng tiếp cận khả thi là theo dõi
thư mục \texttt{oob-logs/} liên tục trong suốt quá trình fuzzing, kể
cả sau khi request gốc đã kết thúc. Token được lưu vào
\texttt{\_token\_map} với timestamp; nếu callback đến muộn, fuzzer
có thể tra cứu token và liên kết với ứng viên gốc theo cơ chế
\textit{deferred lookup} đã đặt nền trong thiết kế hiện tại.

\textbf{Mở rộng sang các lớp lỗ hổng OOB khác.} Phương pháp tín hiệu
OOB không chỉ áp dụng cho SSRF. Các lớp lỗ hổng khác có đặc điểm
tương tự --- không để lại dấu vết trong response --- cũng có thể hưởng
lợi: XXE với DNS exfiltration (\texttt{<!DOCTYPE x [<!ENTITY xxe
SYSTEM "http://oob-server/...">]>}), SQL injection out-of-band
(\texttt{LOAD\_FILE} hay DNS lookup trong MySQL), hay Log4Shell-style
JNDI injection. Kiến trúc OOB server và cơ chế token đã xây dựng có
thể được tái sử dụng trực tiếp.

\textbf{Tăng cường bộ payload bypass.} Tập 21 biến thể hiện tại được
xây dựng dựa trên các kỹ thuật bypass đã biết. Trong thực tế, mỗi
framework và thư viện HTTP có hành vi URL parsing khác nhau, tạo ra
các bypass mới liên tục. Hướng phát triển là tích hợp cơ sở dữ liệu
bypass từ các nguồn như PayloadsAllTheThings và cập nhật bộ mutator
tự động.

\textbf{Tích hợp vào pipeline CI/CD.} Hiện tại hệ thống được chạy
thủ công theo từng target. Một hướng thực tế là đóng gói thành một
Docker image tự chứa có thể tích hợp trực tiếp vào pipeline CI/CD,
tự động quét SSRF mỗi khi có pull request chạm tới code xử lý URL
người dùng.

\textbf{Báo cáo tự động với Proof-of-Concept.} Hiện tại \texttt{SSRFVulnCheck}
ghi report dạng JSON thô. Một bước cải tiến là tự động tạo PoC request
có thể tái hiện lỗ hổng (curl command với payload đúng), cùng với phân
loại mức độ nghiêm trọng dựa trên sink function và mức xác thực yêu
cầu, giúp triage nhanh cho security engineer.
